% Generated draft: requires mathematical and source review.
\documentclass[11pt]{article}
\usepackage[margin=25mm]{geometry}
\usepackage{fontspec}
\setmainfont{Latin Modern Roman}
\usepackage{amsmath,amssymb,mathtools,unicode-math}
\setmathfont{Latin Modern Math}
\usepackage{enumitem,needspace,xurl,hyperref}
\setlength{\parindent}{0pt}
\setlength{\parskip}{5pt}
\setlength{\emergencystretch}{3em}
\newcommand{\N}{\mathbb{N}}
\newcommand{\Z}{\mathbb{Z}}
\newcommand{\Q}{\mathbb{Q}}
\newcommand{\R}{\mathbb{R}}
\newcommand{\C}{\mathbb{C}}
\newcommand{\F}{\mathbb{F}}
\newcommand{\A}{\mathbb{A}}
\newcommand{\PP}{\mathbb{P}}
\newcommand{\E}{\mathbb{E}}
\newcommand{\eps}{\varepsilon}
\newcommand{\dd}{\,\mathrm{d}}
\begin{document}
\section*{UnsolvedMath ID: 30001876}
\textbf{Source problem code:} catalogue entry
\subsection{Definitions and mathematical statement}
A Maass form is understood here as a smooth automorphic eigenfunction of the relevant invariant differential operators, with the automorphy, growth, and normalization conditions still to be specified. Let \(F\) be the \(\operatorname{GL}_3\) Maass form in the source question, and write \(M_3\) for its unspecified automorphic quotient. Let
\[
\iota:N\longrightarrow M_3
\]
denote the codimension-\(2\) submanifold inclusion described in the source as arising from the ``upper-triangular embedding'' of \(\operatorname{GL}_2\). Its restriction is the function
\[
F_N(y)=F(\iota(y)),\qquad y\in N.
\]
The source does not specify the embedding, the quotients, or how central and determinant variables are treated; these symbols therefore record the requested geometry without fixing it.

Let \(u\) be a given \(\operatorname{SL}_2(\Z)\) Maass form. In the usual scalar-valued modular-surface interpretation, its domain is
\[
X=\operatorname{SL}_2(\Z)\backslash\mathbb H,
\qquad \mathbb H=\{x+iy:x\in\R,\ y>0\},
\]
where \(\operatorname{SL}_2(\Z)\) acts by fractional linear transformations. The source does not specify an identification of \(N\) with \(X\), or a procedure for obtaining a comparison function \(u_N\) on \(N\) from \(u\).

\textbf{Source question, with unresolved qualifications explicit.} For a \(\operatorname{GL}_3\) Maass form \(F\) and a given \(\operatorname{SL}_2(\Z)\) Maass form \(u\), determine their restriction inner product. Once a common domain \(N\), a comparison function \(u_N\), and a measure \(\mu_N\) have been specified, the intended pairing would have the form
\[
\langle F_N,u_N\rangle
=\int_N F(\iota(y))\,\overline{u_N(y)}\,\mathrm d\mu_N(y),
\]
provided that this integral exists. More generally, determine the spectral decomposition of the restricted form \(F_N\) on the intended restriction space.

The integral is a conditional formulation of the requested quantity, not a supplied evaluation or convergence assertion. A precise spectral-decomposition problem requires the restriction space, measure, spectral operators, and function space to be fixed, together with any necessary treatment of non-square-integrable terms. The source does not specify whether an explicit coefficient formula, an identity involving other invariants, or another characterization is sought. It supplies no theorem-level assertion resolving these choices.
\subsection{Short English statement}
The question asks how a \(\operatorname{GL}_3\) Maass form decomposes after restriction to a codimension-\(2\) submanifold associated with \(\operatorname{GL}_2\), including its pairing with a given \(\operatorname{SL}_2(\Z)\) Maass form. The embedding, common domain, measure, and spectral setting require clarification before this becomes a fully specified mathematical problem. This is a statement-only draft requiring human mathematical review.
\subsection{Sources}
\begin{itemize}
\item Local UnsolvedMath registry: \texttt{data/unsolvedmath/problems.json}, with \texttt{catalog/problems.json} as the supplementary registry; supplied permanent ID \(30001876\), source problem code \texttt{catalogue entry}, title ``Restriction Spectra of \(\operatorname{GL}_3\) Maass Forms.'' The supplied record belongs to the Oberwolfach Reports Open Problems collection.
\item Supplied report citation: \emph{The Analytic Theory of Automorphic Forms} (2012), Oberwolfach Reports, pp.~2397--2458, DOI \texttt{10.4171/owr/2011/43}. Supplied source URL and literature-source link, labelled ``original OWR report'': \url{https://doi.org/10.4171/owr/2011/43}.
\item Attribution and chronology: the supplied proposed year is 2011; no individual proposer is identified. The report citation gives 2012, while the DOI and supplied workshop-year tag refer to 2011. These supplied dates are retained without resolving their bibliographic relationship.
\item Supplied verification metadata labels the statement \texttt{corrected\_verified} and reports that its mathematical content was checked against the official Oberwolfach report, with notation, context, and question form normalized without changing the source-level claim. This is supplied provenance, not independent verification by this draft.
\item Supplied literature assessment, checked on \texttt{2026-08-21}: status ``open''; the extracted sentence requests a spectral decomposition without specifying a theorem-level claim, and no general resolution is identified. The listed checked source is the original OWR report linked above. This assessment is supplied dated curation, not independently verified; time-sensitive status and cited claims require independent review.
\end{itemize}
\end{document}
